\section{Introduction}

\subsection{Format}

You will be given two 30-minute interviews. The first 25 minutes will aim to reproduce the interview experience and will end with 5 minutes for discussion. The first interview will emulate a tutorial and will cover previously seen material. The second will be on unseen material. If you have any questions, please give them to Mr Hack who will forward them on to the interviewer team. For any online interviews you will be responsible for how to communicate your answers (using Microsoft teams whiteboard and a tablet with stylus for example).

For the first interview, we ask that you attempt \textbf{one} of the questions below and bring your notes to the interview. These are hard questions\footnote{These questions will almost certainly be harder than anything you will be asked to solve on the spot.} so you are not expected to go in the right direction immediately. The purpose is to assess your problem solving creativity, tenacity, and intuition so include your whole thought process. You are encouraged to try more than one question for your own enjoyment, but if you do this then the question you choose to discuss should be the one you believe best demonstrates your problem-solving abilities. You have around one week to work on your problems (note that if this format is used in the actual interviews then you will likely only have one evening).

For the second interview you will be asked an introductory question to ease you into things and the questions will get gradually harder. Do not attempt to judge how well this interview has gone based on how fast you answered the questions.

Full worked solutions to the problems given in this document and those asked during any interviews will be made available afterwards\footnote{This will not happen in your non-mock interviews.}. Hints are often given for each part and it is advised that you attempt problems again using these hints before looking at the full solutions. We ask that you do not share these problems with the year below in case we re-use some of these problems in later years.

\subsection{Interview-Taking Advice}

\begin{itemize}
	\item Relax. Your interview will be between several engineering/physics/maths nerds talking about engineering/physics/maths problems - enjoy it.
	\item Not being able to solve a problem but instead exploring multiple incorrect paths through a problem that ended up being wrong would be a better interview performance than simply writing down the correct answer. The important part is to motivate why you decided on those paths and why you gave up on those paths. This is not a test of whether you can get questions correct - there are entrance exams for this purpose.
	\item Interviewers are also looking for how well they can work with a student in a tutorial and want to hear your thought process, even if you are going down the wrong path or saying something incorrect. This is a natural part of exploring a problem and does not count against you.
	\item Can you work with a tutor? The interview is more collaborative than combative and the interviewer will be trying to help you when they see you are stuck - make sure to use any hints.
	\item You can and should ask clarifying questions\footnote{The interviewer may have deliberately omitted an assumption that changes the question such as whether a function is real-valued, whether friction is present et cetera. You may want to discuss the suitability of making a given assumption and then ask the interviewer whether that it ok once you have decided.}, but avoid asking questions that try to probe the route through the problem - this part is meant to come from you.
	\item If you spot a pattern that ends up being useful, the interviewer may want you to explain/prove that later. It may be helpful to your problem solving process to note that your conjecture is unproven but assume it is true for the time being.
	\item This is also a test of problem solving. Try to demonstrate that you can build intuition for an unseen problem, break a problem down into simpler parts, and apply understanding gained from previously solved problems.
\end{itemize}

\subsection{Problem-Solving Advice}

\begin{itemize}
	\item If the original problem is too hard, consider a simplified version. For example if the problem is very abstract, looking at a concrete example may help.
	\item Make observations about the problem before diving in. For example when deriving the volume of a sphere of radius $r$, you expect it to be proportional to $r^3$. Any simple bounds or relationships you spot can also help.
	\item Graphs are your friend. For example, proving that a cubic having three roots is equivalent to having two stationary points with opposite sign.
	\item Have you solved a similar problem before which may provide useful intuition?
	\item Look for simplifications such as symmetries. For example, when finding the average distance between two randomly chosen points on the circumference of a circle, the circle can always be oriented such that one of the point is at the top of the circle.
	\item If your answer involves considering a large number of cases then you are probably doing something wrong. You should typically only begin listing out cases to try and spot a pattern, but do not intend on finishing.
	\item Considering how much information you need to extract and how much information you can extract is often useful when finding strategies. For example, 3 “yes/no” questions can only distinguish between 8 states.
	\item Understand the pigeonhole principle: \url{https://www.youtube.com/watch?v=TCZ3YwbcDaw}.
	\item Understand how possibilities combine. If $A$ has $n$ possibilities and $B$ has $m$ possibilities then $A$ and $B$ together have at most $nm$ possibilities. Ask your teacher for a probability exercise comparing the configuration space for independent and non-independent events to understand further.
\end{itemize}